\subsection{معرفی و انگیزه}
این مکانیزم بر مبنای مبانی نظری پیشرفته در ادبیات توجه و فشرده‌سازی \cite{kim2022learned} توسعه یافته است.

مکانیزم توجه ترنسفورمر کارآمد با هرس پیشرونده توکن‌ها (\lr{EAT: Efficient Attention Transformer with Progressive Token Pruning}) یکی از راهکارهای پیشرفته در جهت کاهش توأم پیچیدگی زمانی و حافظه‌ای در مدل‌های زبانی است. در کاربردهای پردازش زبان طبیعی، طول دنباله ورودی به ویژه در لایه‌های میانی و انتهایی حاوی تعداد زیادی توکن افزونه یا با بار معنایی ناچیز است؛ با این وجود، توجه کامل استاندارد منابع محاسباتی یکسانی را برای تمامی توکن‌ها تا آخرین لایه صرف می‌کند.

ماژول \lr{eat} با تلفیق مکانیزم توجه چندسره بهینه‌شده، انطباق‌پذیری پارامتر-کارای لورا (\lr{LoRA}) و ماژول هرس پیشرونده توکن‌ها (\lr{ProgressiveTokenPruning})، راهکاری دومنظوره را پیشنهاد می‌دهد. این ماژول امکان انتقال بی‌درنگ وزن‌های پیش‌آموزش‌دیده (\lr{Zero-Shot Weight Inheritance}) را بر روی طیف وسیعی از مدل‌های خانواده DistilBERT، RoBERTa، GPT-2 و Qwen فراهم ساخته و با حذف تدریجی توکن‌های با ضریب اهمیت پایین در طول لایه‌ها، سرعت استنتاج را به شکل چشمگیری افزایش می‌دهد.
